% !TEX root = ../my-thesis.tex
\chapter{Einleitung}
\label{sec:intro}

Wissenschaftliches Rechnen besteht häufig nicht aus einer einzelnen
Berechnung, sondern aus mehreren aufeinander aufbauenden
Rechenschritten. Ausgeführt werden solche Abläufe sowohl auf einzelnen Rechnern als auch auf High-Performance-Computing-Systemen (HPC-Systemen). Mit zunehmender Zahl der Schritte wird es aufwendiger, ihre Ausführung und die dabei entstehenden Daten zu verwalten.

\section{Motivation}
\label{sec:intro:motivation}
Wissenschaftliche Workflows bestehen aus mehreren miteinander
verbundenen Rechenschritten, zwischen denen Daten weitergegeben
werden und die in ihrer Ausführung voneinander abhängen können.
Bei komplexeren Abläufen wird es dadurch zunehmend aufwendig, die
einzelnen Schritte von Hand in der richtigen Reihenfolge
auszuführen und die benötigten Daten zwischen ihnen weiterzugeben.
Workflow-Management-Systeme automatisieren diese Abläufe und
übernehmen dabei unter anderem die Planung und Koordination der
Ausführung, die Übertragung von Daten sowie die Überwachung des
Ausführungszustands~\cite[S.~28\,f.]{badia2017workflows}\cite[S.~2]{suter2026terminology}.
Es steht eine große Zahl solcher Systeme zur Verfügung. Die Wahl
eines Systems wird dabei unter anderem durch die Anforderungen der
jeweiligen Fachgemeinschaft, die Verbreitung einzelner Systeme
oder historische Gründe beeinflusst~\cite[S.~28]{badia2017workflows}. Eine einzelne, für alle wissenschaftlichen Prozesse und Ausführungsumgebungen geeignete Lösung gibt es nicht~\cite[S.~2]{suter2026terminology}. Wer ein System auswählen will, steht damit vor der Frage, wie sich die Systeme in der
Beschreibung von Workflows, im Laufzeitverhalten und im Betrieb
unterscheiden. Für den Einsatz auf einem HPC-System kommt zusätzlich die Frage hinzu, wie sich die Systeme dort bereitstellen und mit dem
Ressourcenverwalter betreiben lassen, der die verfügbaren
Rechenressourcen auf die auszuführenden Aufgaben verteilt.

Hier setzt die vorliegende Arbeit an. Sie vergleicht mit Nextflow,
StreamFlow, Merlin, Pegasus und AiiDA fünf
Workflow-Management-Systeme anhand zweier Workflow-Muster, einer
Pipeline und eines Scatter-Gather-Musters, die in allen Systemen
mit derselben Rechenlogik umgesetzt werden. Die Systeme werden
zunächst auf einem einzelnen Rechner und anschließend auf dem
HPC-System MOGON unter Slurm untersucht. Dabei werden sowohl
qualitative Unterschiede zwischen den Systemen als auch ihr
Laufzeitverhalten betrachtet.


\section{Zielsetzung und Forschungsfragen}
\label{sec:intro:goal}
Ziel der Arbeit ist ein systematischer Vergleich ausgewählter
Workflow-Management-Systeme hinsichtlich ihrer Modellierung,
ihres Laufzeitverhaltens und ihres Einsatzes in einer
HPC-Umgebung. Die Arbeit orientiert sich an vier
Forschungsfragen. Die erste Frage betrifft die Modellierung:
Wie unterscheiden sich die betrachteten
Workflow-Management-Systeme bei der Umsetzung eines Pipeline- und
eines Scatter-Gather-Workflows (F1)? Die zweite Frage richtet sich
auf das Laufzeitverhalten: Welches Laufzeitverhalten zeigen die Systeme
bei der Ausführung dieser Workflows, gemessen an Gesamtlaufzeit (Makespan), Overhead und Koordinationsgrößen (F2)? Die Koordinationsgrößen erfassen dabei
zeitliche Abstände beim Übergang zwischen aufeinanderfolgenden
Schritten und beim Start paralleler Verarbeitungsinstanzen. Die
dritte Frage gilt dem Einsatz auf dem HPC-System: Wie lassen sich
die Systeme auf einem Slurm-basierten HPC-System betreiben, und
wie setzen sich die Laufzeiten unter Clusterbedingungen zusammen
(F3)? Die vierte Frage zielt auf die Einordnung: Für welche Arten
von Workflows und Anwendungsfällen erscheinen die einzelnen
Systeme besonders geeignet (F4)?

\section{Aufbau der Arbeit}
\label{sec:intro:structure}
Die Arbeit ist wie folgt aufgebaut.
Kapitel~\ref{sec:basics} erläutert die Grundlagen wissenschaftlicher
Workflows und HPC-Systeme, stellt die betrachteten
Workflow-Management-Systeme vor und ordnet die Arbeit gegenüber
verwandten Arbeiten ein. Kapitel~\ref{sec:methodik} beschreibt das
Evaluationsdesign, die verwendeten Workflow-Muster und das Vorgehen
für den Vergleich der Systeme. Kapitel~\ref{sec:lokal} beschreibt die Umsetzung der Workflow-Muster in den fünf Systemen, vergleicht sie anhand der qualitativen Kriterien und stellt die Ergebnisse der lokalen Messungen dar.
Kapitel~\ref{sec:mogon} betrachtet anschließend den Einsatz der
Systeme auf dem HPC-System MOGON und die dort durchgeführten
Messungen. Kapitel~\ref{sec:diskussion} diskutiert die Ergebnisse
entlang der Forschungsfragen und ordnet die Eignung der Systeme
sowie die Grenzen der Arbeit ein. Kapitel~\ref{sec:conclusion}
fasst die Ergebnisse zusammen und gibt einen Ausblick.
